% ---
% titre: Reconnaître une situation de proportionnalité (tableaux et situations)
% theme: FA
% chapitre: Proportionnalité
% sous_chapitre: Reconnaître une situation
% niveau: [10e]
% type: EVAL
% tags: [c-proportionnalite,p-question-TS]
% difficulte: 1
% corrige: true
% ---

\question
Pour chaque situation, indique par une phrase s'il s'agit d'une situation proportionnelle ou non.

\begin{parts}
\vspace{20pt}\part[1] Pour préparer 1 pizza, Marc utilise 200\,g de farine. Pour 2 pizzas, il en utilise 400\,g et pour 3 pizzas, 600\,g. La quantité de farine est-elle proportionnelle au nombre de pizzas ? 

\vspace{10pt}\begin{tcolorbox}[enhanced jigsaw, interior hidden, colframe=box, parbox=false]%
Réponse : 
\sol{\vspace{20pt}}{Oui, la situation est proportionnelle. \textbf{(1\,pt)}}
\end{tcolorbox}

\vspace{20pt}\part[1] Il fait 10\,\textdegree C à 8\,h du matin et 18\,\textdegree C à 16\,h. La température est-elle proportionnelle à l'heure de la journée ? 

\vspace{10pt}\begin{tcolorbox}[enhanced jigsaw, interior hidden, colframe=box, parbox=false]%
Réponse : 
\sol{\vspace{20pt}}{Non, la température n'est pas proportionnelle à l'heure de la journée. \textbf{(1\,pt)}}
\end{tcolorbox}

\vspace{20pt}\part[1]   Sarah a 10 ans et pèse 35\,kg. Son papa a 40 ans et pèse 78\,kg. Le poids d'une personne est-il proportionnel à son âge ? 

\vspace{10pt}\begin{tcolorbox}[enhanced jigsaw, interior hidden, colframe=box, parbox=false]%
Réponse : 
\sol{\vspace{20pt}}{Non, le poids d'une personne n'est pas proportionnel à son âge. \textbf{(1\,pt)}}
\end{tcolorbox}

\vspace{20pt}\part[1]    Un ticket de bus coûte 2,50\,CHF. Le prix à payer est-il proportionnel au nombre de tickets achetés ?

\vspace{10pt}\begin{tcolorbox}[enhanced jigsaw, interior hidden, colframe=box, parbox=false]%
Réponse : 
\sol{\vspace{20pt}}{Oui, le prix à payer est proportionnel au nombre de tickets achetés (à moins d'un rabais de groupe). \textbf{(1\,pt)}}
\end{tcolorbox}
\end{parts}
 